\documentclass[11pt]{amsart}

\usepackage[T1]{fontenc}
\usepackage{amsmath,amssymb,amsthm,mathtools}
\usepackage[margin=1in]{geometry}
\usepackage{enumitem}

\newtheorem{theorem}{Theorem}[section]
\newtheorem{proposition}[theorem]{Proposition}
\newtheorem{lemma}[theorem]{Lemma}
\theoremstyle{definition}
\newtheorem{definition}[theorem]{Definition}
\theoremstyle{remark}
\newtheorem{remark}[theorem]{Remark}

\newcommand{\R}{\mathbb R}
\newcommand{\C}{\mathbb C}
\newcommand{\Q}{\mathbb Q}
\newcommand{\Z}{\mathbb Z}
\newcommand{\dd}{\mathrm d}
\newcommand{\PD}{\operatorname{PD}}
\newcommand{\supp}{\operatorname{supp}}
\newcommand{\im}{\operatorname{im}}
\newcommand{\Span}{\operatorname{span}}
\newcommand{\id}{\operatorname{id}}
\newcommand{\End}{\operatorname{End}}

\title{Holomorphic defining coordinates for symplectic normal-crossings divisors}
\date{}

\begin{document}

\begin{abstract}
The answer is affirmative, with the cohomologous deformation allowed in
the question.  The proof first enlarges the divisor by auxiliary
symplectic hypersurfaces so that their cohomology classes satisfy a
redundancy condition.  An integral matrix of relations among these
classes then produces globally defined, smooth, nonvanishing functions
on the divisor complement.  Locally, suitable roots of monomials in
these functions are defining coordinates.  A convexity argument for
tame almost complex structures makes the required functions
holomorphic simultaneously.
\end{abstract}

\maketitle

\section{Statement and conventions}

A \emph{symplectic simple normal-crossings divisor} will mean a finite
union
\[
D=\bigcup_{i=1}^{m}D_i
\]
of closed embedded symplectic submanifolds of real codimension two,
such that every intersection
\[
D_I:=\bigcap_{i\in I}D_i
\]
is transverse and symplectic, and its symplectic orientation agrees
with its intersection orientation.  Empty intersections are permitted.
A labelled hypersurface may be disconnected; this convention does not
change the notion, since its connected components are disjoint.

A smooth complex coordinate system means a smooth real coordinate
system with its coordinates grouped into complex-valued functions.
It is not understood to be holomorphic unless this is explicitly
stated.

\begin{theorem}\label{thm:main}
Let $(M^{2n},\omega)$ be a closed symplectic manifold, and let
$D\subset M$ be a symplectic simple normal-crossings divisor.
There exist a smooth path of symplectic forms $\omega_t$, $0\leq t\leq1$,
and an $\omega_1$-tame almost complex structure $J$ such that:
\begin{enumerate}[label=\textup{(\roman*)}]
\item $\omega_0=\omega$ and $[\omega_t]=[\omega]$ for every $t$;
\item $D$ remains a symplectic simple normal-crossings divisor for
      every $\omega_t$;
\item every point of $M$ has smooth complex coordinates
      $(z_1,\ldots,z_n)$ in which
      \[
      D=\{z_1\cdots z_r=0\}
      \]
      for some $0\leq r\leq n$, and
      \[
      \dd z_j\circ J=i\,\dd z_j,\qquad 1\leq j\leq r.
      \]
\end{enumerate}
\end{theorem}

We use two established results from symplectic geometry:
regularization of symplectic normal-crossings divisors, and relative
approximately holomorphic hypersurface constructions.  We state
precisely the consequences needed below, including an explanation of
the relative and nonintegral versions of the latter.  The remaining
argument is elementary.

\section{Geometric preliminaries}

\subsection{Regularization and compatible almost complex structures}

\begin{lemma}\label{lem:regularization}
Let $E$ be a symplectic simple normal-crossings divisor in a closed
symplectic manifold $(X,\Omega)$.  There is a path
\[
\Omega_t=\Omega+\dd\mu_t,\qquad \mu_0=0,
\]
through symplectic forms for which $E$ remains a symplectic simple
normal-crossings divisor, such that $E$ admits a symplectic
regularization for $\Omega_1$.

Moreover, an $\Omega_1$-compatible almost complex structure $K$ can
be chosen to preserve the tangent bundle of every component of $E$.
\end{lemma}

\begin{proof}
The existence of the exact deformation and the regularization is
\cite[Theorem~2.13]{TMZ}, with the parameter space reduced to a point.
Here and below the numbering refers to the version specified in the
bibliography.
% Source verification for the exact deformation theorem: ([math.stonybrook.edu](https://www.math.stonybrook.edu/~azinger/research/SympDivConf.pdf))

We explain the last assertion.  A regularization supplies, near a
point lying on the components indexed by $I$, a split normal-bundle
model
\[
N=\bigoplus_{i\in I}N_i\longrightarrow E_I,
\]
with the components given by the hyperplane subbundles $v_i=0$.
The symplectic form in this model is
\[
\pi^*\Omega_I+\frac12\sum_{i\in I}\dd(\rho_i\alpha_i),
\]
where $\rho_i$ is a squared Hermitian norm and $\alpha_i$ is the
angular form of a unitary connection; see \cite[Section~2.2]{TMZ}.
In particular, along $v_i=0$, the vertical $i$th line is the
symplectic orthogonal complement of the tangent bundle of that
hyperplane.
% Source verification for the split normal-bundle model: ([math.stonybrook.edu](https://www.math.stonybrook.edu/~azinger/research/SympDivConf.pdf))

Using the connection splitting, choose a local Riemannian metric
making the horizontal space and the individual vertical lines
mutually orthogonal.  Along each component this metric makes
$TE_i$ orthogonal to $(TE_i)^{\Omega_1}$.  Cover $X$ by such
neighborhoods and by neighborhoods disjoint from $E$.  A partition
of unity gives a global metric $h$ with the same orthogonality
property: that property is a linear condition on a metric at a
point of $E_i$.

Define $A\in\End(TX)$ by
\[
h(Au,v)=\Omega_1(u,v).
\]
Then $A$ is invertible and $h$-skew-adjoint.  Set
\[
P=(-A^2)^{1/2},\qquad K=AP^{-1}.
\]
The positive square root is smooth, and $A$ commutes with $P$.
Consequently $K^2=-\id$, $K$ is $h$-orthogonal, and
\[
\Omega_1(u,Ku)=h(Pu,u)>0,\qquad
\Omega_1(Ku,Kv)=\Omega_1(u,v).
\]
Thus $K$ is $\Omega_1$-compatible.

At a point of $E_i$, the orthogonality property of $h$ implies that
$A$ preserves $TE_i$: if $u\in TE_i$ and
$v\in(TE_i)^{\Omega_1}$, then
\[
h(Au,v)=\Omega_1(u,v)=0.
\]
Functional calculus now shows that $P$ and $K$ preserve $TE_i$
as well.
\end{proof}

\subsection{Divisor line bundles}

We use the standard classification of smooth complex line bundles by
their first Chern class, the equality of the first Chern class with
the Euler class of the underlying oriented plane bundle, and the
Chern--Weil formula
\[
c_1(L)_{\R}=\left[\frac{i}{2\pi}F_\nabla\right].
\]
References are \cite[Sections~9 and~14]{MS} and
\cite[Chapter~XII]{KN}.  In particular, a complex line bundle with
vanishing integral first Chern class is smoothly trivial, not merely
trivial after tensoring its cohomology with $\R$.

\begin{lemma}\label{lem:divisor-bundle}
Let $Y\subset X$ be a closed embedded codimension-two submanifold
with a complex structure on its normal bundle $N_Y$.  There are a
smooth complex line bundle $L(Y)\to X$ and a transverse section
$\sigma_Y$ such that
\[
\sigma_Y^{-1}(0)=Y,\qquad c_1(L(Y))=\PD[Y].
\]
The intrinsic derivative along its zero set identifies
\[
\dd\sigma_Y:TX|_Y/TY\longrightarrow L(Y)|_Y
\]
with the identity of $N_Y$.

In particular, if an almost complex structure $K$ preserves $TY$
and the normal complex structure is induced by $K$, then
\[
\dd\sigma_Y\circ K=i\,\dd\sigma_Y
\quad\text{along }Y.
\]
\end{lemma}

\begin{proof}
Identify a tubular neighborhood $T$ of $Y$ with a disk neighborhood
of the zero section of $N_Y$.  On $T$ take the bundle $\pi^*N_Y$
and its tautological section $v\mapsto v$.  On $X\setminus Y$
take the trivial line bundle with section $1$.  Over
$T\setminus Y$, glue these bundles by identifying $1$ with the
nonzero vector $v$.  The sections glue to $\sigma_Y$.
Its derivative on the normal bundle is the identity.

The zero locus of a transverse section represents the Euler class
of the underlying oriented plane bundle, hence the first Chern
class of this complex line bundle.  This gives the asserted
Poincar\'e duality formula, with the orientation determined by
the specified normal complex structure.
\end{proof}

\subsection{The relative hypersurface input}

For a sequence of tensor fields, the notation $O_{C^\infty}(1)$
means boundedness in every fixed derivative norm, with respect to
fixed geometric data on the base.

\begin{lemma}[Relative approximately holomorphic hypersurfaces]
\label{lem:relative-AH}
Let $(X,\Omega,K)$ be a closed symplectic manifold with $K$
$\Omega$-compatible.  Let $Y_1,\ldots,Y_q$ be a finite collection
of closed $K$-invariant submanifolds, including $X$ itself if
desired.  Let $L_k$ be Hermitian line bundles with unitary
connections satisfying
\[
F_{L_k}=-ik\Omega+O_{C^\infty}(1).
\]
For every sufficiently large $k$, there is a smooth section $s_k$
of $L_k$ which is transverse to zero on $X$ and whose restriction
to every $Y_a$ is transverse to zero.  All the resulting zero
sets in $X$ and in the $Y_a$ are symplectic, with their
complex-normal zero-set orientations.
\end{lemma}

\begin{proof}
We recall the precise approximately holomorphic input and its
relative application.  Put $g=\Omega(\,\cdot\,,K\,\cdot\,)$ and
$g_k=kg$.  The curvature hypothesis makes $L_k$ an
asymptotically very ample sequence in the sense of
\cite[Definition~2.1]{Auroux01}: its curvature is positive of
order $k$, its $(0,2)$-part is bounded, and its curvature
derivatives are $O(k)$.

Estimated transversality gives approximately holomorphic sections,
with arbitrarily small approximately holomorphic perturbations
available, for which the restrictions in question are uniformly
transverse.  The relevant general construction is
\cite[Theorem~1.1 and Section~4]{Auroux01}; the
restriction-to-a-prescribed-submanifold version is the main
relative transversality theorem of \cite{Mohsen}.
% Source verification for general asymptotic ampleness and estimated transversality: ([people.math.harvard.edu](https://people.math.harvard.edu/~auroux/papers/pp-seoul.pdf))
% Source verification for transversality of restrictions to a prescribed submanifold: ([numdam.org](https://www.numdam.org/item/AFST_2019_6_28_4_655_0/))

For clarity, the passage to the particular formulation above
requires no extension theorem for holomorphic sections.
Use the ambient Gaussian peak sections centered at points of
$Y_a$, supplied by \cite[Lemma~2.3]{Auroux01}.  Their restrictions
to the $K$-invariant submanifold $Y_a$ have the same local
approximately holomorphic estimates and nonvanishing lower
bounds.  Thus the local perturbations for transversality on
$Y_a$ are restrictions of ambient perturbations.  The
Gaussian-decay globalization argument applies with centers on
$Y_a$.  Apply this successively to the finite list of submanifolds,
taking each perturbation sufficiently small to preserve the
previous uniform transversality bounds.  The openness and
successive-perturbation argument is the one used in
\cite[proof of Theorem~2]{Auroux97}.
% Source verification for the ambient perturbation and openness argument: ([people.math.harvard.edu](https://people.math.harvard.edu/~auroux/papers/ahsub4.pdf))

More explicitly, the resulting sections satisfy, for constants
independent of $k$,
\[
\|s_k\|_{C^2(g_k)}\leq C,\qquad
\|\bar\partial_Ks_k\|_{C^1(g_k)}\leq Ck^{-1/2},
\]
and their derivatives, restricted to each $TY_a$, have right
inverses of uniformly bounded norm wherever the restricted
section is sufficiently small.  This also explains why the
construction applies to the above general sequence $L_k$, rather
than only to tensor powers of one line bundle.

At a zero in $Y_a$, the derivative of $s_k|_{Y_a}$ is therefore
a uniformly surjective real-linear map whose antilinear part
tends to zero.  Its kernel is close to the kernel of a surjective
complex-linear map.  For large $k$ this kernel is symplectic.
The homotopy which removes the small antilinear part stays
surjective, so its zero-set orientation is the symplectic
orientation.  The same argument applies on $X$.  For a
zero-dimensional $Y_a$, transversality means that the section
has no zero there.
\end{proof}

\begin{remark}\label{rem:relative-SNC}
Suppose $E$ is a symplectic simple normal-crossings divisor and
$K$ is compatible with $\Omega$ and preserves all its components.
Apply Lemma~\ref{lem:relative-AH} to all the nonempty intersections
$E_I$, including $E_\varnothing=X$.  The new zero set $H$
then enlarges $E$ to a symplectic simple normal-crossings divisor.

Indeed, transversality on $E_I$ is exactly the additional
transversality required for $E_I\cap H$.  The orientation assertion
in the lemma gives the required agreement of intersection and
symplectic orientations.
\end{remark}

\section{A cohomologically redundant enlargement}

The purpose of the auxiliary hypersurfaces is purely
cohomological.  Importantly, their construction does not require
the original symplectic class to be rational.

\begin{lemma}\label{lem:outside-subspace}
Suppose $E$ is a symplectic simple normal-crossings divisor in
$(X,\Omega)$ and there is an $\Omega$-compatible almost complex
structure preserving its components.  Let
\[
W\subsetneq H^2(X;\Q)
\]
be a proper rational vector subspace.  There is a closed embedded
symplectic hypersurface $H$ such that $E\cup H$ is a symplectic
simple normal-crossings divisor and
\[
\PD[H]\notin W.
\]
\end{lemma}

\begin{proof}
Choose integral classes $\xi_1,\ldots,\xi_b$ whose images form a
basis of $H^2(X;\Z)$ modulo torsion.  Choose closed real forms
$\beta_j$ representing $\xi_j$, and write
\[
[\Omega]=2\pi\sum_{j=1}^{b}\lambda_j\xi_j
\quad\text{in }H^2(X;\R).
\]
There is a real one-form $\eta$ such that
\[
\Omega=2\pi\sum_{j=1}^{b}\lambda_j\beta_j+\dd\eta.
\]

Choose Hermitian line bundles $B_j$ with $c_1(B_j)=\xi_j$ and
unitary connections of curvature $-2\pi i\beta_j$.
Such connections exist: start with any unitary connection and
correct it by an imaginary-valued one-form, using the
Chern--Weil formula.

Let $m_j(k)$ be an integer nearest to $k\lambda_j$, and set
\[
L_k=\bigotimes_{j=1}^{b}B_j^{\otimes m_j(k)}.
\]
Equip $L_k$ with the tensor-product connection, modified by
the imaginary-valued one-form $-ik\eta$.  Its curvature is
\begin{align*}
F_{L_k}
&=-2\pi i\sum_j m_j(k)\beta_j-ik\,\dd\eta\\
&=-ik\Omega
  -2\pi i\sum_j\bigl(m_j(k)-k\lambda_j\bigr)\beta_j\\
&=-ik\Omega+O_{C^\infty}(1).
\end{align*}
Write $h_k=c_1(L_k)$.

Choose an integral class $a$ whose rational image is not in $W$,
and a Hermitian line bundle $A$ with $c_1(A)=a$, equipped with
any unitary connection.  The sequence $L_k\otimes A$ satisfies
the same curvature estimate.

By Lemma~\ref{lem:relative-AH} and
Remark~\ref{rem:relative-SNC}, for all sufficiently large $k$
both classes
\[
h_k,\qquad h_k+a
\]
are represented by hypersurfaces which can separately be added
to $E$ as symplectic simple normal-crossings components.
They cannot both belong to $W$, since their difference is $a$.
Choose the hypersurface corresponding to a class outside $W$.
\end{proof}

\begin{lemma}\label{lem:redundant-enlargement}
Starting with $(M^{2n},\omega,D)$, one can make a finite sequence
of the exact deformations allowed in Theorem~\ref{thm:main} and
add finitely many auxiliary hypersurfaces to obtain a symplectic
simple normal-crossings divisor
\[
E=\bigcup_{i=1}^{N}E_i,\qquad D\subset E,
\]
in $(M,\Omega)$ with the following properties:
\begin{enumerate}[label=\textup{(\alph*)}]
\item there is an $\Omega$-compatible almost complex structure
      preserving every $E_i$;
\item for every $I\subset\{1,\ldots,N\}$ with $|I|\leq n$,
      the classes
      \[
      \{\PD[E_j]:j\notin I\}
      \]
      span $H^2(M;\Q)$.
\end{enumerate}
Throughout the deformations, the original divisor $D$ is fixed
as a subset and remains symplectic simple normal crossings.
\end{lemma}

\begin{proof}
Set $b=\dim H^2(M;\Q)$.  First apply
Lemma~\ref{lem:regularization} to the current divisor.
We construct $n+1$ disjoint groups of auxiliary labels, each
group containing $b$ hypersurfaces whose classes form a basis
of $H^2(M;\Q)$.

Within one group, suppose $\ell<b$ linearly independent classes
have already been obtained.  Apply
Lemma~\ref{lem:outside-subspace} with $W$ their span.  Add the
resulting hypersurface to the current divisor.  Apply
Lemma~\ref{lem:regularization} to this enlarged divisor before
proceeding to the next step.  Its deformation is exact and
preserves all the components already constructed.  This
procedure constructs the required basis group, and repeating
it constructs all $n+1$ groups.

At the end, the compatible almost complex structure from the
last regularization gives (a).  Deleting at most $n$ labels can
meet at most $n$ of the $n+1$ groups, so at least one complete
basis group remains.  This proves (b).  If $b=0$, the spanning
condition is vacuous and no auxiliary components are necessary.

All forms encountered have cohomology class $[\omega]$.
The finitely many deformation paths can be concatenated
smoothly by making their parametrizations constant near their
endpoints.
\end{proof}

\subsection{An integral relation matrix}

\begin{lemma}\label{lem:matrix}
Let $c_1,\ldots,c_N\in H^2(M;\Z)$.  Suppose that, after deleting
any set of at most $n$ indices, the remaining rational classes
still span the rational span of all the $c_i$.
There is a symmetric integer matrix
\[
B=(b_{ij})_{1\leq i,j\leq N}
\]
such that
\[
\sum_{j=1}^{N}b_{ij}c_j=0
\quad\text{in }H^2(M;\Z),\qquad 1\leq i\leq N,
\]
and every nonempty principal submatrix
\[
B_I=(b_{ij})_{i,j\in I},\qquad |I|\leq n,
\]
is positive definite.
\end{lemma}

\begin{proof}
Choose a rational basis of the span of the classes and let $C$
be the matrix whose $j$th column is the coordinate vector of
$c_j$.  Thus $C$ has full row rank.  The orthogonal projection
onto its kernel is the rational symmetric matrix
\[
P=I_N-C^{\mathsf T}(CC^{\mathsf T})^{-1}C.
\]
If the rational span is zero, use $P=I_N$.

We claim that $P_I$ is positive definite when $|I|\leq n$.
If $v\in\R^N$ is supported on $I$ and
$v^{\mathsf T}Pv=0$, then $Pv=0$, since $P$ is an orthogonal
projection.  Hence $v=C^{\mathsf T}u$ for some $u$.
For $j\notin I$ the $j$th entry of $v$ is zero, so $u$ is
orthogonal to every column of $C$ outside $I$.  Those columns
span the whole column space, by hypothesis.  Thus $u=0$ and
$v=0$, proving the claim.

Multiply $P$ by a positive integer to clear its denominators.
Every resulting row gives an integral combination of the $c_j$
which vanishes rationally, hence is torsion.  The torsion
subgroup of $H^2(M;\Z)$ is finite.  Multiplying once more by
its exponent kills all these torsion classes.  The resulting
integer matrix $B$ has the asserted properties.
\end{proof}

\section{Gluing tame structures while preserving holomorphic functions}

The following observation is the gluing mechanism in the proof.
It is useful that holomorphicity of a fixed function becomes an
affine-linear condition in suitable coordinates on the space of
tame almost complex structures.

\begin{lemma}\label{lem:tame-gluing}
Let $(X,\Omega)$ be a symplectic manifold.  Suppose $\{U_\alpha\}$
is a locally finite open cover, $J_\alpha$ is an $\Omega$-tame
almost complex structure on $U_\alpha$, and $\{\rho_\alpha\}$
is a subordinate smooth partition of unity.

There is a smooth $\Omega$-tame almost complex structure $J$
with the following pointwise property.  If
$\lambda:T_xX\to\C$ is a complex-valued real-linear covector and
\[
\lambda\circ J_\alpha(x)=i\lambda
\]
for every $\alpha$ with $\rho_\alpha(x)>0$, then
\[
\lambda\circ J(x)=i\lambda.
\]
\end{lemma}

\begin{proof}
Choose an $\Omega$-compatible reference structure $J_*$.
Write
\[
TX\otimes_{\R}\C=V_-\oplus V_+,
\]
where $V_-$ and $V_+$ are the $-i$ and $+i$ eigenbundles of
$J_*$.  On the complexified tangent bundle consider the
Hermitian form
\[
H(v,w)=i\Omega_{\C}(v,\overline w).
\]
It is positive definite on $V_-$, negative definite on $V_+$,
and these two subbundles are $H$-orthogonal.

For any $\Omega$-tame almost complex structure $J'$, its
$-i$ eigenspace is $H$-positive.  Indeed, its vectors have the
form $u+iJ'u$, and
\[
H(u+iJ'u,u+iJ'u)=2\Omega(u,J'u)>0
\]
when $u\neq0$.  Projection to $V_-$ therefore identifies this
eigenspace with the graph of a complex-linear map
\[
A':V_-\longrightarrow V_+,\qquad \|A'\|<1,
\]
where the norms come from $H$ on $V_-$ and $-H$ on $V_+$.

Conversely, the graph of any such contraction is $H$-positive.
Its conjugate is $H$-negative, so the two graphs are
complementary.  They define a real almost complex structure,
which is tame by the same calculation.

Let $A_\alpha$ represent $J_\alpha$ in these graph coordinates,
and set
\[
A=\sum_\alpha\rho_\alpha A_\alpha.
\]
The products are extended by zero outside their respective
domains.  This is smooth and satisfies
\[
\|A\|\leq\sum_\alpha\rho_\alpha\|A_\alpha\|<1.
\]
Let $J$ be the tame structure defined by its graph.

The condition $\lambda\circ J_\alpha=i\lambda$ is equivalent
to the complex-linear extension of $\lambda$ vanishing on the
$-i$ eigenspace, or
\[
\lambda(v)+\lambda(A_\alpha v)=0
\qquad(v\in V_-).
\]
This is affine-linear in $A_\alpha$.  Taking the weighted sum
proves the corresponding identity for $A$, and hence the
claimed property of $J$.
\end{proof}

\section{Monomial functions and holomorphic normal coordinates}

\begin{proposition}\label{prop:redundant-case}
Let $E=\bigcup_{i=1}^{N}E_i$ be a symplectic simple
normal-crossings divisor in a closed symplectic manifold
$(X^{2n},\Omega)$.  Assume:
\begin{enumerate}[label=\textup{(\alph*)}]
\item there is an $\Omega$-compatible almost complex structure
      $K$ preserving every $E_i$;
\item the classes $c_i=\PD[E_i]$ satisfy the hypothesis of
      Lemma~\ref{lem:matrix}.
\end{enumerate}
Then there is an $\Omega$-tame almost complex structure $J$
for which, near every point, all the local branches of $E$
have simultaneous $J$-holomorphic defining coordinates.
\end{proposition}

\begin{proof}
We divide the argument into four steps.

\medskip
\noindent\textbf{Step 1: global functions on the complement.}
For each $E_i$, use the normal complex structure induced by $K$
and Lemma~\ref{lem:divisor-bundle} to obtain a complex line bundle
$L_i$ and a section $\sigma_i$ satisfying
\[
\sigma_i^{-1}(0)=E_i,\qquad
c_1(L_i)=c_i,\qquad
\dd\sigma_i\circ K=i\,\dd\sigma_i\quad\text{along }E_i.
\]
Choose the matrix $B$ from Lemma~\ref{lem:matrix}.  Each line
bundle
\[
P_i:=\bigotimes_{j=1}^{N}L_j^{\otimes b_{ij}}
\]
has zero integral first Chern class, and is therefore smoothly
trivial.  Choose a complex-linear trivialization
\[
\tau_i:P_i\longrightarrow X\times\C.
\]
A negative tensor power denotes a power of the dual line.

On $X\setminus E$, all the sections $\sigma_j$ are nonzero.
They therefore define smooth functions
\[
f_i
=
\tau_i\!\left(\bigotimes_{j=1}^{N}\sigma_j^{\otimes b_{ij}}\right)
:
X\setminus E\longrightarrow\C^*.
\]
For negative exponents, the inverse section of the dual line
is used.  No holomorphicity is asserted yet.

\medskip
\noindent\textbf{Step 2: local simultaneous monomialization.}
Fix $p\in E$, and put
\[
I=I(p):=\{i:p\in E_i\}.
\]
Then $1\leq |I|\leq n$.  Take a sufficiently small neighborhood
$U$ of $p$ which meets no component $E_j$ with $j\notin I$,
and trivialize all the $L_j$ there.  Write
\[
\sigma_j=g_j e_j
\]
in these frames.  For $j\in I$, put $z_j=g_j$.
Transversality of the divisor means that the map
$(z_j)_{j\in I}$ is a submersion near $p$.  Moreover,
\[
\dd z_j(p)\circ K(p)=i\,\dd z_j(p),\qquad j\in I.
\]

Complete these functions to smooth complex coordinates
\[
\bigl((z_j)_{j\in I},y_1,\ldots,y_{n-|I|}\bigr)
\]
whose differential at $p$ is complex-linear for $K(p)$.
To do this, extend the independent complex covectors
$\dd z_j(p)$ to a complex basis of the cotangent space, choose
smooth functions with these additional differentials, and
apply the inverse function theorem.

For $i\in I$, the definition of $f_i$ gives
\[
f_i=a_i\prod_{j\in I}z_j^{b_{ij}}
\quad\text{on }U\setminus E,
\]
where each $a_i:U\to\C^*$ is smooth: all factors belonging to
indices outside $I$ are nonvanishing and are absorbed into
$a_i$.

Shrink $U$ so that $a_i$ admits a smooth logarithm $h_i$ for
each $i\in I$.  Since $B_I$ is invertible, define
\[
u_j=\sum_{i\in I}(B_I^{-1})_{ji}h_i,\qquad
w_j=z_j e^{u_j},\qquad j\in I.
\]
Then
\begin{equation}\label{eq:monomial}
f_i=\prod_{j\in I}w_j^{b_{ij}},
\qquad i\in I,
\quad\text{on }U\setminus E.
\end{equation}
Also,
\[
\{w_j=0\}=E_j\cap U,\qquad
\dd w_j(p)=e^{u_j(p)}\dd z_j(p).
\]
Thus, after shrinking once more,
\[
\bigl((w_j)_{j\in I},y_1,\ldots,y_{n-|I|}\bigr)
\]
is a smooth complex coordinate system whose differential at
$p$ is complex-linear for $K(p)$.

Let $J_p$ be the pullback of the standard complex structure
in these coordinates.  It agrees with $K$ at $p$, so it is
$\Omega$-tame on a sufficiently small neighborhood $U_p$.
Every $f_i$ with $i\in I(p)$ is $J_p$-holomorphic on
$U_p\setminus E$, by \eqref{eq:monomial}.

At a point outside $E$, choose any local $\Omega$-tame almost
complex structure and give the corresponding neighborhood the
empty label set.  We have therefore constructed an open cover
by neighborhoods $U_\alpha$, with label sets $I_\alpha$, such
that:
\begin{enumerate}[label=\textup{(\arabic*)}]
\item $U_\alpha\cap E_i=\varnothing$ if $i\notin I_\alpha$;
\item $J_\alpha$ is $\Omega$-tame on $U_\alpha$;
\item $f_i$ is $J_\alpha$-holomorphic on $U_\alpha\setminus E$
      whenever $i\in I_\alpha$.
\end{enumerate}
Take a finite such cover, using compactness of $X$.

\medskip
\noindent\textbf{Step 3: gluing.}
Choose a smooth partition of unity $\{\rho_\alpha\}$
subordinate to this finite cover, with
$\supp\rho_\alpha\subset U_\alpha$.  Apply
Lemma~\ref{lem:tame-gluing} to obtain an $\Omega$-tame almost
complex structure $J$.

For each $i$, define the open neighborhood
\[
V_i
=
X\setminus
\bigcup_{\alpha:\,i\notin I_\alpha}\supp\rho_\alpha.
\]
Property (1) above shows that $E_i\subset V_i$.
If $x\in V_i\setminus E$, every index with
$\rho_\alpha(x)>0$ has $i\in I_\alpha$.  Thus every contributing
$J_\alpha$ satisfies
\[
\dd f_i(x)\circ J_\alpha(x)=i\,\dd f_i(x).
\]
The gluing lemma gives
\begin{equation}\label{eq:f-holomorphic}
\dd f_i\circ J=i\,\dd f_i
\quad\text{on }V_i\setminus E.
\end{equation}

\medskip
\noindent\textbf{Step 4: recovery of the defining coordinates.}
Return to $p\in E$, set $I=I(p)$, and use the local coordinates
constructed in Step~2.  Restrict their neighborhood to
\[
U\cap\bigcap_{i\in I}V_i.
\]
On its complement of $E$, all the functions $f_i$ and $w_j$
appearing in \eqref{eq:monomial} are nonzero.  Logarithmic
differentiation yields
\[
\frac{\dd f_i}{f_i}
=
\sum_{j\in I}b_{ij}\frac{\dd w_j}{w_j},
\qquad i\in I.
\]
By \eqref{eq:f-holomorphic}, the covectors on the left are
complex-linear for $J$.  Since $B_I$ is invertible, the same
is true of every $\dd w_j/w_j$, and hence of every $\dd w_j$:
\[
\dd w_j\circ J=i\,\dd w_j
\quad\text{on }U\setminus E.
\]
Both sides are smooth on all of $U$.  The complement of the
finite union of codimension-two submanifolds $E$ is dense.
Consequently this identity holds throughout $U$.

The functions $w_j$, $j\in I$, are therefore simultaneous
$J$-holomorphic defining coordinates, and
\[
E\cap U=\left\{\prod_{j\in I}w_j=0\right\}.
\]
This proves the proposition.
\end{proof}

\section{Proof of the main theorem}

\begin{proof}[Proof of Theorem~\ref{thm:main}]
If $D$ is empty, choose any $\omega$-compatible almost complex
structure and take $r=0$ everywhere.  The zero-dimensional
case is likewise immediate.  Assume henceforth that $n\geq1$
and $D\neq\varnothing$.

Apply Lemma~\ref{lem:redundant-enlargement}.  It produces an
endpoint symplectic form $\Omega$, connected to $\omega$ by a
smooth path of cohomologous symplectic forms preserving the
original divisor, and an enlarged divisor
\[
E=D\cup E_{\mathrm{aux}}
\]
satisfying the hypotheses of
Proposition~\ref{prop:redundant-case}.  Let $J$ be the
$\Omega$-tame almost complex structure provided by that
proposition.

At a point $p\in D$, the proposition gives smooth complex
coordinates in which every branch of the enlarged divisor
through $p$ is a coordinate hyperplane, with its defining
coordinate $J$-holomorphic.  Retain as the first
$z_1,\ldots,z_r$ precisely the coordinates corresponding to
the original components of $D$ through $p$.  Use the remaining
coordinates, including any auxiliary defining coordinates, to
complete the smooth complex coordinate system.  Shrinking
the neighborhood to avoid all other original components gives
\[
D=\{z_1\cdots z_r=0\},\qquad
\dd z_j\circ J=i\,\dd z_j\quad(1\leq j\leq r).
\]
At a point outside $D$, choose a neighborhood disjoint from
$D$ and take $r=0$.

Finally, set $\omega_1=\Omega$ and use the concatenated path
from Lemma~\ref{lem:redundant-enlargement}.  It has all the
properties asserted in the theorem.
\end{proof}

\begin{remark}
The argument does not confuse holomorphic defining functions
with the weaker condition that the Nijenhuis tensor have
tangential image along the divisor.  Holomorphicity here holds
on actual open neighborhoods.  It follows from the global
functions $f_i$, the gluing identity
\eqref{eq:f-holomorphic}, and the invertibility of the principal
matrices $B_I$.
\end{remark}

\begin{remark}
Neither the deformation nor the auxiliary construction changes
the cohomology class of the symplectic form.  The rounding
construction in Lemma~\ref{lem:outside-subspace} is what permits
approximately holomorphic hypersurfaces to be used without
assuming that $[\omega]$ is rational.  Separately, the
torsion-killing step in Lemma~\ref{lem:matrix} ensures that the
tensor-product line bundles used to define the functions
$f_i$ are genuinely trivial.
\end{remark}

\begin{thebibliography}{99}

\bibitem{Auroux97}
D.~Auroux,
\emph{Asymptotically holomorphic families of symplectic submanifolds},
Geom. Funct. Anal. \textbf{7} (1997), 971--995.

\bibitem{Auroux01}
D.~Auroux,
\emph{Estimated transversality in symplectic geometry and projective maps},
in \emph{Symplectic Geometry and Mirror Symmetry}
(Seoul, 2000), World Scientific, Singapore, 2001, pp.~1--30;
arXiv:math/0010052.

\bibitem{TMZ}
M.~Farajzadeh Tehrani, M.~McLean, and A.~Zinger,
\emph{Normal crossings singularities for symplectic topology},
arXiv:1410.0609.
The references to Theorem~2.13 and Section~2.2 use the authors'
version dated April~30, 2017.

\bibitem{KN}
S.~Kobayashi and K.~Nomizu,
\emph{Foundations of Differential Geometry}, Vol.~II,
Wiley, New York, 1969.

\bibitem{MS}
J.~W.~Milnor and J.~D.~Stasheff,
\emph{Characteristic Classes},
Annals of Mathematics Studies, vol.~76,
Princeton University Press, Princeton, NJ, 1974.

\bibitem{Mohsen}
J.-P.~Mohsen,
\emph{Transversalit\'e quantitative en g\'eom\'etrie symplectique:
sous-vari\'et\'es et hypersurfaces},
Ann. Fac. Sci. Toulouse Math. (6) \textbf{28} (2019), no.~4.

\end{thebibliography}

\end{document}